\chapter{Example: numerical convergence}\label{ch:ex-numerics}

A Floquet calculation of the junction involves six independent approximations: the Floquet
cutoff $M$, the static basis size $K$ (if used), the QP cutoff, the bath discretization, the
number of drive samples and, for the propagator, the number of time steps. This chapter
studies the Floquet cutoff, the static basis and the bath for the phase-driven reference
junction at QP cutoff $\q{numerics.cutoff}$; convergence with the QP cutoff must be checked as for the
static problem (\cref{sec:qp}).
The sampling and the time steps are checked by the midpoint test of \cref{sec:sampling} and
the step-doubling estimate of \cref{sec:propagator}.

\section{The Floquet cutoff}

\Cref{fig:numerics}(a) shows the error of the tracked $S$ and $S^*$ quasienergies for a
strong phase drive, $A=\q{numerics.strong_amplitude}$ at $\Omega=\q{numerics.omega}$, as a
function of $M$, together with the rigorous bound \eqref{eq:residual-bound}. The error
decreases exponentially, by one to two orders of magnitude per harmonic, because the harmonics
of the drive decay like Bessel functions.
The bound always lies above the error, which is roughly the square of the bound, as expected
for an isolated eigenvalue. The bound is conservative and can be used as a stopping criterion. The
weight in the outermost blocks, \code{boundary_weights}, decreases at the same rate.

\section{The static basis}

In the static basis the many-body space is truncated to the $K$ lowest undriven eigenstates.
\Cref{fig:numerics}(b) compares the tracked quasienergies with those of the full Fock space
at QP cutoff $\q{numerics.cutoff}$ (dimension $\q{numerics.fock_dimension}$) for a weaker
drive.
The error is about $\q{numerics.small_basis_error}$ already for $K\approx10$ and then
decreases only slowly; it falls below $\q{numerics.large_basis_error}$ when $K$ approaches the
full dimension. The reason is physical. The phase drive couples the subgap states to
quasiparticle excitations at all energies up to the bandwidth $D=\q{numerics.bandwidth}$, and
each contributes a small AC Stark shift \eqref{eq:ac-stark}. The static basis treats
couplings among the retained states exactly; recovering the accumulated virtual
contributions requires sufficiently many off-resonant states.
\code{fock_residuals} provides a conservative error bound. For quantitative AC Stark shifts use the
Fock basis; for resonances, splittings and scans the static basis is efficient.

\section{The discretized continuum}

\Cref{fig:numerics}(c,d) follows the dressed singlet across the one-photon threshold
$\Omega=E_D+\Delta-E_S\approx\q{numerics.threshold}$, above which its first replica lies in
the continuum, for baths with $\q{numerics.bath_pairs}$ level pairs per lead. Below the
threshold the AC Stark shifts obtained with the $\q{numerics.bath_count}$ baths agree within
$\q{numerics.below_spread}\%$, and the calculation is converged in the bath. The exceptions are
a few narrow avoided crossings at two-photon resonances, $2\Omega=E_a-E_S$, with discrete
continuum levels $a$; like the one-photon crossings above the threshold, they move with the
bath. Above the threshold, the dressed singlet shows a narrow avoided crossing, with a dip of
its static weight to about one half, wherever a replica of a discrete continuum level crosses
it. The positions of these crossings differ from bath to bath: they are discretization
artefacts (\cref{sec:continuum}). In the continuum limit they merge into a finite lifetime of
the dressed bound state, set by the rate at which the drive excites quasiparticles out of it.

\begin{figure}[htb]
\centering
\manualfigure{numerics}
\caption{Convergence for the phase-driven reference junction. (a) Error of the tracked $S$
and $S^*$ quasienergies versus the Floquet cutoff (static basis,
$K=\q{numerics.cutoff_basis_states}$, $A=\q{numerics.strong_amplitude}$,
$\Omega=\q{numerics.omega}$), with the rigorous bound \eqref{eq:residual-bound}. (b) Error of
the static-basis result versus $K$, relative to the full Fock space
($A=\q{numerics.weak_amplitude}$), with the Fock-space residual. (c) AC Stark shift of the
dressed singlet and (d) its static weight versus $\Omega$ across the one-photon continuum
threshold (dotted), for $\q{numerics.bath_count}$ bath sizes ($A=\q{numerics.weak_amplitude}$);
the legend gives the number of level pairs per lead and the Fock-space dimension.}
\label{fig:numerics}
\end{figure}

\section{A practical workflow}

\begin{enumerate}
\item Converge the static problem first: QP cutoff and bath, for the states of interest
(\cref{ch:static}).
\item Choose $M$ from the modulation index: $M\gtrsim(\text{largest level modulation})/\Omega+5$;
check that \code{residuals} and \code{boundary_weights} are small.
\item For scans, use the static basis and check \code{fock_residuals}, or compare selected
points with Fock-space shift-invert results.
\item Cross-check one point with \code{floquet_propagator}, which has no Floquet cutoff.
\item Above continuum thresholds, compare at least two bath sizes and interpret only
features that are stable under refinement.
\end{enumerate}

\tutorial{05_convergence}
